\chapter{Design under observation}
\label{ch:design}

\section*{What this chapter is for}

Between discovery and building there is a drawing, and it is watched. The
drawing is not assessed on being correct \dash{} it is assessed on whether you
can think in front of somebody, build something incrementally, and stop.

\section{One shape, one sentence}

The technique that survives observation is to draw in phases, where each phase
is one new shape and one sentence saying why it is there. Never a finished
diagram appearing in silence, and never a shape drawn without being named.

\begin{kitmethod}
A phase is:

\textbf{One shape} added to the picture \dash{} a box, an arrow, a boundary,
a label.

\textbf{One sentence} saying what it is and why it exists.

\textbf{A pause}, short, in which somebody can interrupt.

Ten to fifteen phases is a whole architecture. At roughly thirty to forty
seconds each, that is under ten minutes.
\end{kitmethod}

The pause is doing the work. An interviewer who interrupts is telling you
something is wrong or interesting, and that is the cheapest correction
available. \judgement{A candidate who draws for six minutes in silence and then
says ``any questions?'' has converted a conversation into a presentation, and
presentations are not what this round is for.}

\section{A phase order that generalises}

Not a template to memorise \dash{} an order that works because each phase is
answerable from the one before.

\begin{kittable}{@{}llX@{}}
\textbf{\#} & \textbf{Shape} & \textbf{The sentence} \\
\midrule
1 & the actor & who this is for \\
2 & the thing they want & the one path that must work \\
3 & the boundary & what is inside this system and what is not \\
4 & the store & what has to survive a restart \\
5 & the entities & the three or four nouns, and their relationships \\
6 & the second actor & the one discovery found \\
7 & the operation that can conflict & where two people collide \\
8 & the external dependency & what you do not control \\
9 & what happens when 8 fails & the degradation \\
10 & the seam & where the next thing would attach \\
\end{kittable}

Phase 7 is the one that reads as senior, because it is the phase most
candidates never reach. Every system has an operation where two actors can
collide over one resource, naming it unprompted signals that you have operated
something, and it gives you a concrete thing to build in the slice that follows.

\section{Two scenes}

Your screen is doing two jobs and they need different layouts. Decide in
advance and rehearse the switch, because doing it live costs a minute and looks
like fumbling.

\textbf{The design scene:} the drawing surface large, the written artefact from
discovery visible beside it, nothing else. No editor, no terminal, no
notifications.

\textbf{The build scene:} editor and terminal, with the design still reachable
in one keystroke \dash{} you will return to it, and returning to it deliberately
is a good moment.

\judgement{Whatever you draw in cannot matter as much as being fluent in it.
The failure mode is not an ugly diagram; it is four minutes lost to a tool you
use twice a year, at the exact moment you are being watched.}

\section{What to draw and what to skip}

\begin{kittable}{@{}XX@{}}
\textbf{Draw} & \textbf{Do not draw} \\
\midrule
The boundary of your system & every component inside it \\
What survives a restart & the deployment topology \\
The entities and their relationships & every field \\
The one operation that can conflict & the happy path in sequence-diagram detail \\
What you do not control & a class diagram \\
\end{kittable}

The rule underneath: draw what constrains the build. Anything that does not
change what you write in the next hour is decoration, and decoration under a
clock is a decision to spend time on the wrong thing in front of an audience.

\section{The design is a claim, not a monument}

You will change it. Say so while drawing \dash{} name the part you are least
sure of, and what would make you change it. Then, when the build reveals it was
wrong, changing it is the process working rather than a retreat.

\begin{kitnote}
``I am least confident about this boundary; if it turns out that the second
actor needs to see this state, it moves'' is a strong sentence at minute
twenty-five, and an even stronger one at minute seventy when it turns out the
second actor did need to see that state.
\end{kitnote}

\begin{drill}
Draw a system you know well in ten phases, out loud, timed, alone. Record it.
Watch it back with one question: could somebody have interrupted you? If the
phases run together and there is nowhere to break in, that is the defect, and
it is invisible from the inside.
\end{drill}

\section*{What this chapter does not claim}

The phase order is the author's, assembled from what an interviewer needs in
order to follow a design they have not seen. It is not measured, and a
different order that reaches phase 7 is as good.

It makes no claim that the design will be assessed for correctness. Where the
criteria are published they are about thinking and judgement, and this chapter
is written for that.
